The continuing emphasis on miniaturization and functionality in materials design has brought one-dimensional (1D) nanostructures to the forefront of nanotechnology research. Out of these 1D nanostructures, single walled carbon nanotubes (SWCNTs) - cylindrical allotropes of carbon with unusual mechanical, electronic, and thermal properties - remain the model materials for the building blocks of next generation devices. The special aspects of SWCNTs, including chirality-dependent metallicity/semiconductivity, high-mobility carrier transport, and tunability of optoelectronic properties, put SWCNTs in a unique position to study hybrid structures based on 1D systems. These hybrids, or 1D heterostructures, are a mix of one-component functional molecules, or materials, with SWCNTs to produce materials with new properties.

This document outlines a discovery internship at the Charles Coulomb Laboratory, undertaken under the direction of Dr. Laurent Alvarez, that examined the optical spectroscopy of 1D heterostructures, notably Raman spectroscopy. The internship involved studying how the encapsulation of organic molecules (4T:quaterthiophene) into SWCNTs can modify their vibrational, electronic and structural parameters. Raman spectroscopy (a non-destructive, and non-invasive laser technique sensitive to lattice vibrations) is an effective way of characterising these molecules:

\begin{itemize}
    \item Chirality and diameter via Radial Breathing Modes (RBMs),
\end{itemize}

\begin{itemize}
    \item Electronic transitions through resonance effects,
\end{itemize}

\begin{itemize}
    \item Defect density (D-band),
\end{itemize}

\begin{itemize}
    \item Strain, doping, and interfacial interactions (G-band and 2D-band shifts).
\end{itemize}

The internship gave me a hands-on experience of experimental methods from SWCNT synthesis, constructing heterostructures, and advanced Raman measuring at different excitation wavelengths (532 nm, 633 nm, 785 nm). The internship had the following goals:


\begin{itemize}
    \item Understanding the synthesis and intrinsic properties of SWCNTs and their chiral vector (n,m)-dependent electronic structure.
\end{itemize}


\begin{itemize}
    \item Using resonant Raman spectroscopy to probe the vibrational response of pristine SWCNTs and heterostructures encapsulated by 4T (4T@HIPCO).
\end{itemize}


\begin{itemize}
    \item Interpreting spectral energy shifts (G-band energy downshifts, RBM energy shifts) associated with charge transfer, mechanical strain, and thermal changes elicited by molecular encapsulation.
\end{itemize}
This report is structured as follows:

\begin{itemize}
    \item Chapter 2 reviews the fundamentals of carbon nanotubes, including atomic structure, chirality classification, electronic properties, synthesis techniques (arc discharge, CVD, laser ablation), and phonon dispersion.
\end{itemize}

\begin{itemize}
    \item Chapter 3 details Raman spectroscopy principles (non-resonant/resonant regimes) and its specific application to SWCNTs.
\end{itemize}

\begin{itemize}
    \item Chapter 4 presents experimental results and discusses the implications of 4T encapsulation on SWCNT properties, supported by multi-wavelength Raman analysis.
\end{itemize}

Through this work, the internship highlighted the importance of optical spectroscopy and the advancement of nanomaterial characterization and mapping the interplay of theoretical placement versus experimental realization into new functionalities in 1D heterostructures.